\subsection{Boundary models for one-sided and mixed extensions}
\label{nb:statements}

We retain the circuit matrices $Q_1,\ldots,Q_6$, their span $L_0$,
and $v(s)=(1,s,s^2)^{\mathsf T}$ from~\eqref{eq:Q} and the
row--column construction. Thus $\sum_{i=1}^6Q_i=0$, and $A_0,B_0$
denote the old three-dimensional coordinate subspaces. All graph and
image closures below have their reduced structures. Normalization is
taken in the variety's own function field.

For nonnegative integers $s,m,k,\alpha,\beta$, define the subalgebras
\begin{align}
 \mathcal R_{s,m}
 &=\C\left[t,\mu_j,y_{ij},\frac{\mu_jy_{ij}}t:
          1\leq j\leq s,\ 1\leq i\leq m\right],
 \label{nb:R}\displaybreak[0]\\
 \mathcal B_{k,\alpha,\beta}
 &=\C\left[t,p,q_j,x_j,z_j,u_h,v_{ij},
      \frac{pz_j+q_jx_j}t,\frac{pu_h}t,\frac{q_jv_{ij}}t:
      1\leq j\leq k,\ 1\leq h\leq\beta,\ 1\leq i\leq\alpha
      \right].\label{nb:B}
\end{align}
Each is the displayed subalgebra of the Laurent polynomial ring in its
algebraically independent base variables with $t$ inverted. Empty index
sets are allowed. In $\mathcal B$, $\alpha$ counts the $v$-rows and
$\beta$ the $u$-variables. Give every base variable degree one; the
displayed fractions also have degree one. The \emph{homogeneous origin}
is the point where all these generators vanish.

We use \emph{smooth equivalence} in the following precise sense: a
smooth local cover of the indicated germ is \'etale over the product
of the model germ with a smooth germ. The covers include the choices
of factor vectors and the ordering of configurations. No canonical
coordinate isomorphism on the unframed normalization is asserted.
The \emph{tangent excess} of a local ring is its embedding dimension
minus its Krull dimension.

\begin{theorem}[One-sided extensions, D]\label{nb:D}
Let
\[
 r\geq6,\qquad a\geq r-2,\qquad b\geq3,\qquad c\geq r,
 \qquad a+b\geq r+2.
\]
In $A=\C^a$, choose distinct new coordinate vectors
$u_1,\ldots,u_{r-5}$ complementary to $A_0$, and put
\begin{equation}\label{nb:row-witness}
 \begin{split}
 R_s&=u_s\otimes v(s+3)\quad(1\leq s\leq r-5),\\
 L_r&=L_0\oplus\langle R_1,\ldots,R_{r-5}\rangle,\\
 T_{\mathrm{row}}
 &=\sum_{i=1}^5Q_i\otimes f_i+
              \sum_{s=1}^{r-5}R_s\otimes f_{5+s}.
 \end{split}
\end{equation}
Here $v(s+3)\in B_0\subset B=\C^b$, even when $b>3$, and
$f_1,\ldots,f_r$ are coordinate vectors in $\C^c$.
The ordinary rank and border rank of $T_{\mathrm{row}}$ equal $r$.
Its entire normalization fibre in $X_r(a,b,c)$ consists of exactly
$r+1$ reduced points. Six correspond to the precisely six actual
unordered $r$-term decompositions and are smooth. The other $r-5$
points correspond to the dependent configurations containing all six
$Q$-directions and all $R$-directions except one; they are not actual
$r$-term decompositions.

Put $\ell=r-6$ and $m=b-3$. Every boundary point is smoothly equivalent
to the homogeneous origin of $\Spec\mathcal R_{\ell,m}$. This ring is
normal, with tangent excess $\ell m$; its origin is smooth precisely
when $\ell m=0$. Thus all boundary points are smooth precisely when
$r=6$ or $b=3$, and all are singular when $r\geq7$ and $b\geq4$.
Every such point has dimension $r(a+b+c-2)$ and embedding dimension
\[
 r(a+b+c-2)+(r-6)(b-3).
\]
There is a Zariski open neighbourhood of $T_{\mathrm{row}}$ over which
normalization is unramified. If $r=6$ or $b=3$, its entire source over
a sufficiently small such neighbourhood is smooth. All conclusions
transport under permutation of the tensor factors together with the
witness; in particular, transposition gives the added-column version.
\end{theorem}

The cases $r=6$ and $b=3$ recover the corresponding one-sided families
above. The disjoint diagonal witnesses are different. Theorem~\ref{nb:D}
does not apply to simultaneous row and column additions and makes no
global smoothness, Cohen--Macaulayness, rational-singularity or
exact-depth assertion about the normalization.

\begin{theorem}[One row and several columns, E1]\label{nb:E1}
Let
\[
 r\geq7,\qquad a\geq4,\qquad b\geq r-3,\qquad c\geq r,
 \qquad a+b\geq r+2,
 \qquad k=r-6.
\]
Choose a new $A$ coordinate $u$ and new $B$ coordinates
$w_1,\ldots,w_k$, and set
\begin{equation}\label{nb:E1-witness}
 \begin{split}
 R&=u\otimes v(4),\qquad S_j=v(j+3)\otimes w_j,\\
 L&=L_0\oplus\C R\oplus\langle S_1,\ldots,S_k\rangle,\\
 T&=\sum_{i=1}^5Q_i\otimes f_i+R\otimes f_6+
                         \sum_{j=1}^kS_j\otimes f_{6+j}.
 \end{split}
\end{equation}
This is $T_{r;4,r-3}$, of support dimensions $(4,r-3,r)$, padded to
the stated format. Its rank and border rank equal $r$, and it has
exactly six actual unordered $r$-term decompositions. The entire
normalization fibre consists of exactly $r+1$ reduced points indexed
by the $r$-element subsets of
$\{[Q_1],\ldots,[Q_6],[R],[S_1],\ldots,[S_k]\}$, as follows:
\begin{enumerate}
\item The six subsets omitting a $Q$-point give smooth points and
      the six actual decompositions.
\item The subset omitting $R$ gives one singular point with local
      model $\mathcal R_{r-6,a-3}$ and tangent excess $(r-6)(a-3)$.
\item Each subset omitting an $S_j$ gives a singular point with local
      model $\mathcal B_{r-7,a-4,b-r+4}$ and tangent excess
      $(r-7)(a-3)+b-r+4$.
\end{enumerate}
The latter configurations are dependent, not actual decompositions.
All the displayed excesses are positive, including at $r=7$, when
$\mathcal B_{0,a-4,b-3}$ is $\mathcal R_{1,b-3}$ after renaming
$p,u_h$. Every normalization point has dimension $D=r(a+b+c-2)$;
its embedding dimension is $D$ at a smooth point and $D$ plus the
stated excess at a boundary point. Normalization is unramified over
some Zariski open neighbourhood of $T$. Transposition gives the
family with one added column and $r-6$ added rows, with all formulas
transported by the swap.
\end{theorem}

\begin{theorem}[Two rows and two columns, E2]\label{nb:E2}
Let $a,b\geq5$, $a+b\geq11$, and $c\geq9$. Choose new coordinates
$u_1,u_2$ in $A$ and $w_1,w_2$ in $B$, and put
\begin{equation}\label{nb:E2-witness}
 \begin{split}
 R_1&=u_1\otimes v(4),\qquad R_2=u_2\otimes v(5),\\
 S_1&=v(4)\otimes w_1,\qquad S_2=v(5)\otimes w_2,\\
 L&=L_0\oplus\langle R_1,R_2,S_1,S_2\rangle,\\
 T&=\sum_{i=1}^5Q_i\otimes f_i+R_1\otimes f_6+R_2\otimes f_7
                         +S_1\otimes f_8+S_2\otimes f_9.
 \end{split}
\end{equation}
This is $T_{9;5,5}$, with support dimensions $(5,5,9)$, padded to
the stated format. Its rank and border rank equal nine, and it has
exactly six actual unordered nine-term decompositions. Its entire
normalization fibre consists of ten reduced points indexed by the
nine-element subsets of the ten displayed matrix directions:
\begin{enumerate}
\item Six smooth points omit a $Q$-point and correspond to the six
      actual decompositions.
\item The two points omitting $R_1$ or $R_2$ are singular, with
      model $\mathcal B_{2,a-4,b-5}$ and excess $2a+b-11$.
\item The two points omitting $S_1$ or $S_2$ are singular, with
      model $\mathcal B_{2,b-4,a-5}$ and excess $a+2b-11$.
\end{enumerate}
The four boundary configurations are dependent. Both excesses are
at least five. The cases $b=5$ and $a=5$ are included, giving
$\beta=0$ in the respective star model. Every normalization point
has dimension $D=9(a+b+c-2)$ and embedding dimension $D$ plus its
listed excess, with excess zero at the smooth points. Normalization
is unramified over a Zariski open neighbourhood of $T$.
\end{theorem}

Theorems~\ref{nb:E1} and~\ref{nb:E2} concern exactly the displayed
mixed witnesses and their admissible padding. Their conclusions
transport under factor permutations with the witnesses. They make
no assertion for broader mixed splits with several retained rows and
columns, nor any global regularity, Cohen--Macaulayness, rational-
singularity or exact-depth assertion about a normalization.

\subsection{The configuration graph and the exact dependence ideals}
\label{nb:geometry}

We prove the three theorems together, retaining the exact parameters
in each statement. Write $W=A\otimes B$,
$S=\Seg(\PP A\times\PP B)\subset\PP W$, and let $H$ be the closure
in $\Gr(r,W)$ of the spans of $r$ independent points of $S$.

\begin{lemma}\label{nb:configuration}
For each displayed witness, $\PP(L)\cap S$ is the reduced scheme
of its $r+1$ displayed matrix directions. The only relation is the
six-core relation. Exactly six $r$-subsets are independent; the
remaining subsets have span dimension $r-1$. The unordered span map
is birational onto $H$, and $\dim H=r(a+b-2)$.

The reduced closure $G$ of the independent unordered span graph in
$\operatorname{Hilb}^r(S)\times H$ is integral, proper and birational
over $H$. Its fibre over $L$ is the reduced scheme of all $r+1$
$r$-element subsets. Its six independent-configuration points are
smooth.
\end{lemma}

\begin{proof}
The core computation of~\ref{norm:core} gives the reduced six-point
intersection in $L_0$. The avoidance calculation in the row--column
construction gives, for every integer $z\geq4$,
\begin{equation}\label{nb:avoidance}
 L_0\cap(A_0\otimes\C v(z))=0,
 \qquad L_0\cap(\C v(z)\otimes B_0)=0.
\end{equation}
For the first equality, the determinant of the rows
$\mathbf1,v(z)^{\mathsf T},v(z)^{\mathsf T}/(1,2,3)$ is
$-z(z-1)(z-3)/6$, and the remaining third row eliminates the last
coefficient. For the second, writing $v(z)b^{\mathsf T}\in L_0$
gives $b_j=jd/(j-z)$ and $(j+z)d=je$ for $j=1,2,3$;
the first two equations force $d=0$ and then $b=0$.
The $v$-directions are distinct and avoid the core factor directions.

Let $E_A,E_B$ be the added row and column spaces. Every matrix in
$L$ has zero $E_A\otimes E_B$ block. For a rank-one matrix
$x\otimes y$, this forces $x_{E_A}=0$ or $y_{E_B}=0$. A nonzero
new row forces the latter equality and its assigned $v$-direction;
distinct directions kill all other new rows, and~\eqref{nb:avoidance}
kills the core block. The matrix is that $R$-direction. Transposition
gives the $S$-directions. If neither new block is present, the matrix
lies in $L_0$. Padding introduces no further possibility: applying
the quotient maps shows that a nonzero rank-one matrix in a tensor
product of subspaces has both factors in those subspaces.

At a $Q$-point, a new row or column in the rank-one tangent space
has its core factor direction, whereas its possible component in $L$
has a distinct $v$-direction. All added coefficients vanish, leaving
the core scalar tangent line. At an $R$-point, the zero mixed block
kills the varying new-column component; distinct row directions kill
the other new rows, and~\eqref{nb:avoidance} kills the core component.
Only the scalar tangent remains. The $S$-case is the transpose, and
padding adds separate complementary tangent blocks. The projective
tangent spaces therefore vanish. Nakayama's lemma applied to
$\mathfrak m/\mathfrak m^2=0$ proves reducedness.

The added matrices occupy distinct row or column blocks, so their
only relation is $\sum_iQ_i=0$. The independent subsets are precisely
those omitting a $Q$-point. The third flattening has rank $r$ and
image $L$. An actual $r$-term decomposition must therefore choose a
basis among these rank-one directions, and that basis determines its
third factors uniquely. There are precisely six actual unordered
decompositions; flattening rank also proves that ordinary rank and
border rank are $r$. Their addition differentials are injective
modulo summand scalings: in an independent third-factor basis, every
matrix variation lies both in $L$ and in its rank-one tangent space,
so it is scalar; the corresponding third-factor variation cancels
exactly that scaling.

For generic factor lists and every subset $J$ of size $j\geq2$,
\[
 \rank(a_i:i\in J)+\rank(b_i:i\in J)
 =\min(a,j)+\min(b,j)\geq j+2.
\]
The last inequality holds because $a,b\geq3$ and $a+b\geq r+2$.
Sylvester's inequality makes every combination supported on at least
two terms have matrix rank at least two. Thus the chosen matrices
are independent and their span has exactly its original rank-one
directions. The generic unordered configuration is unique, so in
characteristic zero the span map is birational. Its domain has
dimension $r(a+b-2)$.

The graph closure is integral and proper over $H$ by projectivity of
the Hilbert scheme. The condition $Z\subset\PP(K)$ is closed: it is
the vanishing of the universal evaluation of the annihilator of the
tautological $K$ in the rank-$r$ bundle $H^0(Z,\OO_Z(1))$. It holds
on the dense independent graph and hence on its reduced closure.
Its fibre over $L$ is consequently a closed subscheme of
$\operatorname{Hilb}^r(S\cap\PP(L))$, the discrete reduced scheme
of $r$-subsets. Repeated or infinitesimal configurations cannot occur,
and a pair $(Z,L)$ determines at most one underlying graph point.
The six independent subsets give smooth graph points, since there
the graph is the graph of the span morphism on the smooth
unordered configuration space.

Every dependent subset occurs. If $R_s=u_s\otimes v(s+3)$ is
omitted, write $v(s+3)=\sum_{i=1}^3\lambda_i r_i^{\mathsf T}$,
where $r_i$ are the rows of $M$, and deform
\begin{equation}\label{nb:row-arc}
 Q_i(\epsilon)=(e_i+\epsilon\lambda_i u_s)r_i
 \quad(1\leq i\leq3),
\end{equation}
keeping the other core and retained added matrices fixed. The sum
of the six core matrices is $\epsilon R_s$. A relation for
$\epsilon\ne0$, projected to the core, makes all six core
coefficients equal. The retained added blocks then kill their own
coefficients, and the omitted row kills the common coefficient.
The matrices are independent. Five deformed core matrices, the
retained additions and the fixed $R_s$ form a basis with independent
limits, so their spans tend to $L$.

If $S_s=v(s+3)\otimes w_s$ is omitted, choose $\eta_j$ with
$v(s+3)=\sum_j\eta_j(-m_j)$ and deform instead
\begin{equation}\label{nb:column-arc}
 Q_{3+j}(\epsilon)=(-m_j)(e_j+\epsilon\eta_jw_s)^{\mathsf T}
 \quad(1\leq j\leq3).
\end{equation}
Projection to the old core and then to the retained and omitted
column blocks gives the same independence and limiting-span argument.
These are the arcs of~\ref{norm:row-column}. Every subset is thus
present, proving the asserted reduced fibre. This alone does not
exclude additional points after normalization; that requires the
local normality proved below.
\end{proof}

\begin{lemma}[Exact boundary charts]\label{nb:charts}
At the boundary configurations of Theorems~\ref{nb:D}--\ref{nb:E2},
the graph germs are smoothly equivalent to the homogeneous origins
of the local algebras listed in those theorems.
\end{lemma}

\begin{proof}
Locally order the distinct points, a free finite \'etale cover, and
choose nonzero factor-vector representatives, a smooth torus torsor.
Call the resulting smooth space $F$. Near such a configuration, the
Hilbert-scheme graph is the ordinary graph of the rational span map
on the smooth configuration space.

The $r$ representative columns have rank $r-1$. Fix an invertible
$(r-1)$-column coordinate minor and let $I$ be the ideal of all
$r\times r$ minors. Row elimination identifies $I$ exactly with the
residual coordinates of the final column modulo the first $r-1$.
Introduce $r-1$ relation coefficients, normalizing one core coefficient
to one. The five other core coefficients are units; the coefficients
of the retained additions vanish at the configuration. The pivot
minor solves the auxiliary coefficients uniquely as regular functions
on $F$. If $w_i$ are their differences from the solved values, the
vector-relation ideal on $F\times\mathbb A^{r-1}$ is
\begin{equation}\label{nb:auxiliary}
 J=(w_1,\ldots,w_{r-1},I).
\end{equation}
Absorb the five unit core coefficients into their first factor vectors.
With core matrices $A_c,B_c$ and retained-factor matrices $U,V$, the
relation becomes
\begin{equation}\label{nb:relation}
 A_cB_c^{\mathsf T}+U\Delta V^{\mathsf T}=0,
\end{equation}
where $\Delta$ is diagonal in the retained coefficients. No vanishing
coefficient is inverted.

Split the core columns into triples and rows by the old supports:
\[
 A_c=\begin{pmatrix}A_{11}&A_{12}\\ A_{21}&A_{22}\end{pmatrix},
 \quad B_c=[B_1\ B_2],\quad B_2=\binom{B_{20}}{B_{21}},
 \quad U=\binom{U_0}{U_1},\quad V=\binom{V_0}{V_1}.
\]
The blocks $A_{11},B_{20}$ are invertible at the circuit. On the
neighbourhood where they remain invertible, use the substitutions
\begin{equation}\label{nb:schur}
 \begin{gathered}
 X=A_{11}^{-1}A_{12},\quad Y=A_{21}A_{11}^{-1},\quad
 P=A_{22}-A_{21}X,\quad C_0=B_1+B_2X^{\mathsf T},\\
 U'=U_1-YU_0,\qquad V'=V_1-B_{21}B_{20}^{-1}V_0.
 \end{gathered}
\end{equation}
These are regular invertible coordinate changes: they solve uniquely
for $A_{22},B_1,U_1,V_1$, retaining the other coordinates. The top
three $A$ rows of~\eqref{nb:relation} solve
\[
 C_0^{\mathsf T}=-A_{11}^{-1}U_0\Delta V^{\mathsf T}.
\]
Subtracting $Y$ times those rows leaves
$PB_2^{\mathsf T}+U'\Delta V^{\mathsf T}$. Its first three $B$
columns solve
\[
 P=-U'\Delta V_0^{\mathsf T}(B_{20}^{\mathsf T})^{-1}.
\]
The remaining equations are exactly the entries of
\begin{equation}\label{nb:residual}
 U'\Delta(V')^{\mathsf T}.
\end{equation}
All operations are invertible operations on ideals. The solved
blocks give $d_0=3b+3(a-3)=3a+3b-9$ independent coordinate
generators $L_1,\ldots,L_{d_0}$.

We identify the residual ideal in each case. For D, there are
$\ell=r-6$ retained rows. Their $U'$ columns are independent new
coordinate directions at the witness, so a fixed $\ell\times\ell$
row minor gives a regular left inverse. Multiplication by this inverse
identifies the ideal of~\eqref{nb:residual} with that of
$\Delta(V')^{\mathsf T}$; the reverse containment follows by
multiplication by $U'$. With $\mu_j$ the coefficients and $y_{ij}$
the entries of $V'$, the ideal is
\begin{equation}\label{nb:row-residual}
 (\mu_jy_{ij}:1\leq j\leq r-6,\ 1\leq i\leq b-3).
\end{equation}
For $\ell=0$ it is zero. The $\mu,y$ are independent vanishing
coordinates; the retained $U'$ data are smooth parameters.

For E1 with its single row omitted, all $r-6$ retained additions
are columns. A regular left inverse of the independent $V'$ columns
identifies~\eqref{nb:residual} with $U'\Delta$. Its ideal is
\begin{equation}\label{nb:column-residual}
 (q_jy_{ij}:1\leq j\leq r-6,\ 1\leq i\leq a-3),
\end{equation}
again in independent vanishing coordinates, with the remaining
$V'$ data smooth parameters.

All other cases have one retained row and $k$ retained columns,
after transposition if necessary. At the witness $U'_R\ne0$ and
$V'_R=0$, whereas $U'_{S_j}=0$ and the $V'_{S_j}$ are independent.
Regular invertible operations in $A/A_0$ send $U'_R$ to $e_1$;
independently, operations in $B/B_0$ send $V'_{S_j}$ to $e_j$.
The required coordinate and minor are units, and their data remain
smooth parameters. These changes act invertibly on the opposite
vanishing columns. Write
\[
 \alpha=a-4,\quad\beta=b-3-k,\qquad
 U'_{S_j}=(x_j,v_{1j},\ldots,v_{\alpha j})^{\mathsf T},\quad
 V'_R=(z_1,\ldots,z_k,u_1,\ldots,u_\beta)^{\mathsf T}.
\]
If $p$ and $q_j$ are the row and column coefficients, the complete
nonzero entry list of~\eqref{nb:residual} is
\begin{equation}\label{nb:star-residual}
 pz_j+q_jx_j\ (1\leq j\leq k),\qquad
 pu_h\ (1\leq h\leq\beta),\qquad
 q_jv_{ij}\ (1\leq i\leq\alpha,\ 1\leq j\leq k).
\end{equation}
All other entries vanish identically. In particular there are no
terms $q_jx_l$ for $l\ne j$ and no omitted higher-order remainders.
For E1 with a column omitted,
$(k,\alpha,\beta)=(r-7,a-4,b-r+4)$, with $\beta\geq1$;
the case $k=0$ leaves just $pu_h$. For E2 with a row omitted,
$(k,\alpha,\beta)=(2,a-4,b-5)$, where $\alpha\geq1$ and
$\beta\geq0$; if $\beta=0$, admissibility gives $\alpha\geq2$.
For a column omitted, transpose to obtain $(2,b-4,a-5)$.

Let $\mathcal Q$ denote the relevant quadratic list. We have the
exact ideal $J=(L_1,\ldots,L_{d_0},\mathcal Q)$. We must remove the
auxiliary relation coordinates while preserving this embedded model.
Every $w_i$ in~\eqref{nb:auxiliary} belongs to $J$. All nonlinear
coordinates vanish at the witness, so $d\mathcal Q=0$ there and each
$dw_i$ is a linear combination of the $dL_j$. The $r-1$ differentials
$dw_i$ are independent, since they are graph coordinates for the
auxiliary coefficients. One coefficient minor is therefore a unit.
Express $w_i$ in the displayed ideal generators and replace those
$r-1$ linear coordinates by $w_i$. Inverting that coefficient minor
preserves the generated ideal together with the remaining $L_j$ and
$\mathcal Q$. The Jacobian is invertible while the nonlinear and
unused coordinates are retained. Restricting to $w=0$ returns the
original $F$ and gives the exact \'etale-local ideal
\begin{equation}\label{nb:dependence-ideal}
 I=(l_1,\ldots,l_\kappa,\mathcal Q),\qquad
 \kappa=d_0-(r-1)=3a+3b-r-8.
\end{equation}
The nonlinear variables remain independent. No radical, initial
ideal or higher-order truncation has been taken. The number $\kappa$
counts linear generators and is independent of the smooth-parameter
count in the rank-five local-ring calculation.

The span map is given in Pl\"ucker coordinates by the maximal minors,
so its graph is the blowup $\Proj\bigoplus_{n\geq0}I^n$.
Equivalently, this is the closure of the graph of the projective
rational map defined by these generators. Flat base change commutes
with ideal powers, so the description is preserved by the \'etale
coordinates and the factor-vector cover.

Use the arc~\eqref{nb:row-arc} or~\eqref{nb:column-arc} for the
omitted addition. Choose the pivot from five independent old-core
entries and one detecting entry for each retained addition: its own
new row with an old $B$ coordinate, or its own new column with an
old $A$ coordinate. None uses the omitted row or column. On this arc,
the normalized pivot relation is still the six-core relation and
all retained coefficients are identically zero. The unchanged old-core
entries first determine the six core coefficients; the distinct
retained entries then force all other coefficients to vanish.
Consequently every residual quadratic is identically zero on the
arc, including when an opposite Schur-substituted factor varies to
first order. In D, one can also see this from the fact that every
$B$ factor stays in $B_0$, so $V'=0$ throughout the arc.

The relation residual has nonzero derivative equal to the omitted
matrix modulo the retained span. Since $d\mathcal Q=0$, this gives
a nonzero linear direction in~\eqref{nb:dependence-ideal}. Choose
$t=l_1$ nonzero on it and replace the other $l_i$ by combinations
whose derivative on the arc is zero. The $t$-chart is the actual
subalgebra
\begin{equation}\label{nb:graph-chart}
 \OO_F[l_2/t,\ldots,l_\kappa/t,\mathcal Q/t]
 \ \subset\ \OO_F[t^{-1}].
\end{equation}
Writing $l_i=th_i$ splits off smooth coordinates $h_i$, centered at
zero at the limit, together with the unused parameters. The remaining
algebra is exactly $\mathcal R_{r-6,b-3}$ in D,
$\mathcal R_{r-6,a-3}$ in the omitted-row case of E1, and the stated
$\mathcal B$ in all star cases. Every model base variable vanishes
at the witness. Every nonlinear numerator vanishes along the whole
arc, so each ratio by $t$ has limit zero. The point is the homogeneous
origin in each case, including opposite omitted and retained types.

The subalgebra~\eqref{nb:graph-chart} retains full $t$-saturation.
Equations merely setting $t$ times a fraction equal to its numerator
need not present it. For example, when $\ell=1$ in the monomial
case, the additional two-by-two minors of
$\bigl(\begin{smallmatrix}t&y_{11}&\cdots&y_{m1}\\
\mu_1&\mu_1y_{11}/t&\cdots&\mu_1y_{m1}/t\end{smallmatrix}\bigr)$
are also necessary. The arcs place the selected point in the actual
graph closure, without an unsaturated presentation. This proves the
asserted smooth equivalences.
\end{proof}

\subsection{Normality of the boundary algebras}
\label{nb:algebra}

\begin{proposition}\label{nb:monomial}
For all $s,m\geq0$, $\mathcal R_{s,m}$ is normal. At its homogeneous
origin its dimension is $1+s+sm$, its embedding dimension is
$1+s+2sm$, and it is smooth precisely when $sm=0$.
\end{proposition}

\begin{proof}
Represent exponents by $(K;A_j;B_{ij})$, with $K$ the $t$ exponent.
The exponent monoid is exactly the integer vectors satisfying
\begin{equation}\label{nb:monoid}
 A_j\geq0,\quad B_{ij}\geq0,\qquad
 K+\sum_j\min\{A_j,\sum_iB_{ij}\}\geq0.
\end{equation}
Indeed, if $p_{ij}$ copies of $\mu_jy_{ij}/t$ occur, then
$\sum_i p_{ij}\leq A_j$ and $p_{ij}\leq B_{ij}$, proving necessity.
If $K\geq0$, base variables suffice. If $K<0$, distribute $-K$
integer units among groups $j$ with capacities
$\min\{A_j,\sum_iB_{ij}\}$, and then within each group among the
integer capacities $B_{ij}$. Successive allocation gives integers
$p_{ij}$ with total $-K$ and the required bounds. The remaining
exponents are supplied by base variables. This proves sufficiency.

The final inequality is equivalent to the finite family
\[
 K+\sum_{j\in E}A_j+\sum_{j\notin E,i}B_{ij}\geq0
 \qquad(E\subset\{1,\ldots,s\}).
\]
The generator lattice is the full integer lattice because all base
variables occur. Thus the monoid is its rational cone intersected
with that lattice. More explicitly, intersect the Laurent polynomial
ring in all base variables with the monomial-weight valuation rings
for these inequalities and the nonnegativity inequalities. A weight
valuation of a Laurent polynomial is the minimum of its monomial
weights; it is multiplicative because lowest homogeneous parts
multiply nontrivially in the Laurent polynomial domain. The
intersection consists exactly of Laurent polynomials whose exponents
satisfy~\eqref{nb:monoid}, hence is $\mathcal R_{s,m}$. Each ring in
the intersection is integrally closed, proving normality, including
empty index ranges.

The total-degree functional $K+\sum_jA_j+\sum_{i,j}B_{ij}$ gives
every displayed generator degree one. They are distinct Laurent
monomials, so they have no linear relation and form a basis of the
cotangent space at the origin. There are $1+s+sm$ independent base
variables and $sm$ fractions. The fraction field has transcendence
degree $1+s+sm$, while the cotangent dimension is $1+s+2sm$.
For $sm=0$ the ring is polynomial; otherwise the positive tangent
excess makes the origin singular.
\end{proof}

\begin{lemma}[Regular-parameter deformation]\label{nb:deformation}
Let $D$ be a local domain obtained by localizing a finitely generated
$\C$-domain, and let $0\ne f$ lie in its maximal ideal. If $D/fD$
is a normal domain, then $D$ is normal.
\end{lemma}

\begin{proof}
The normalization $N$ of $D$ is finite, since finite-type
$\C$-domains have finite normalization and normalization commutes
with localization. The prime ideal $(f)$ has height one by the
principal ideal theorem: $D$ is a domain and $f$ is nonzero and
not a unit. Hence $D_{(f)}$ is a one-dimensional local domain with
principal maximal ideal, so it is a discrete valuation ring.
Every element of $N$ is integral over $D_{(f)}$, whence
$N\subset D_{(f)}$. Localizing $N$ at $D\setminus(f)$ gives
$D_{(f)}$ itself. There is exactly one prime $Q$ of $N$ above
$(f)$, with $N_Q=D_{(f)}$, and $f$ has valuation one at $Q$.

Every height-one prime of $N$ containing $f$ contracts to $(f)$.
Indeed, a finite birational extension of affine $\C$-domains
preserves prime heights: total dimensions agree and the corresponding
residue fields have the same transcendence degree, since their
extension is finite. The dimension formula gives equality of heights,
which persists under the localizations defining $D,N$. A height-one
contraction containing $(f)$ must equal $(f)$. Thus $Q$ is the only
height-one prime of $N$ containing $f$.

Reduction in $N_Q=D_{(f)}$ embeds $N/Q$ in $\Frac(D/fD)$.
It is integral over the normal domain $D/fD$, so $N/Q=D/fD$.
Given $n\in N$, choose $d\in D$ with the same reduction modulo $Q$.
The fraction $(n-d)/f$ has nonnegative valuation at $Q$, because the
numerator has valuation at least one. At every other height-one
prime, $f$ is a unit, so the valuation is again nonnegative. A
Noetherian normal domain is the intersection of its height-one
valuation rings in its fraction field; hence $(n-d)/f\in N$.
Therefore $N=D+fN$. The finite $D$-module $N/D$ satisfies
$N/D=f(N/D)$, and Nakayama's lemma gives $N/D=0$.
\end{proof}

\begin{corollary}\label{nb:graded-deformation}
Let $E$ be a finitely generated nonnegatively graded $\C$-domain,
and let $\xi$ be a nonzero homogeneous element of positive degree.
If $E/(\xi)$ is a normal domain, then $E$ is normal.
\end{corollary}

\begin{proof}
Lemma~\ref{nb:deformation} proves normality at every prime containing
$\xi$. The nonnormal locus is closed and invariant under the
multiplicative action of the grading. Choose finitely many homogeneous
algebra generators, all of nonnegative degree. Every closed point
has an affine limit under scalar multiplication as the scalar tends
to zero: positive-degree coordinates tend to zero and degree-zero
coordinates remain fixed. The limit lies on $\xi=0$. A nonempty
nonnormal locus would contain a closed point, its nonzero-scalar
orbit and its limit, contradicting normality on that fibre. The
nonnegative grading is used here; normality along one fibre of an
arbitrary nonproper family would not suffice.
\end{proof}

\begin{proposition}[Star algebra]\label{nb:star}
For every $k,\alpha,\beta\geq0$, the exact subalgebra
$\mathcal B_{k,\alpha,\beta}$ of~\eqref{nb:B} is normal. Its
dimension and the embedding dimension at its homogeneous origin are
\[
 d=2+3k+\beta+\alpha k,\qquad
 N=2+4k+2\beta+2\alpha k.
\]
The tangent excess is $k(\alpha+1)+\beta$. The origin is smooth
precisely when $k=\beta=0$.
\end{proposition}

\begin{proof}
Write $R=\mathcal B_{k,\alpha,\beta}$ and abbreviate its fractions by
\[
 G_j=\frac{pz_j+q_jx_j}t,\qquad P_h=\frac{pu_h}t,
 \qquad Q_{ij}=\frac{q_jv_{ij}}t.
\]
The additional elements
\begin{equation}\label{nb:H}
 H_{ij}=\frac{pz_jv_{ij}}t=v_{ij}G_j-x_jQ_{ij}
\end{equation}
belong to $R$. Let
$B=\C[t,p,q,x,z,u,v,P,Q,H]$, a monomial algebra; then
$R=B[G_1,\ldots,G_k]$. Give $q_j$ degree one and all other free
variables degree zero. Thus $Q_{ij}$ has degree one, $P_h,H_{ij}$
degree zero, and the leading part of $G_j$ is
$E_j=q_jx_j/t$. We prove
\begin{equation}\label{nb:star-graded}
 \operatorname{gr}R=A:=B[E_1,\ldots,E_k]
\end{equation}
for this total-$q$ filtration.

For $S_0\subset\{1,\ldots,k\}$, put
$B_{S_0}=B[E_l:l\in S_0]$ and fix $j\notin S_0$. The only
fractional generators consuming $q_j$ in this monomial algebra are
$Q_{ij}$, and $x_j$ occurs only as a free variable. These facts
remain true at every intermediate stage. For each $n\geq0$,
\begin{equation}\label{nb:denominator-ideals}
 \{b\in B_{S_0}:bE_j^n\in B_{S_0}\}
       =(t,v_{1j},\ldots,v_{\alpha j})^n.
\end{equation}
To prove this, first let $b$ be a Laurent monomial and choose a
generator representation of $bE_j^n$. Let $s$ be its number of
$Q_{ij}$ factors with this fixed $j$, counted with multiplicity.
If $b$ has $q_j,x_j$ exponents $a_q,a_x$, the representation has
exactly $a_q+n-s$ free $q_j$ factors and $a_x+n$ free $x_j$ factors;
both $a_q,a_x$ are nonnegative. If $s\geq n$, remove $n$ such
$Q$ factors and $n$ free $x_j$ factors. Their product is
$q_j^nx_j^n/t^n$ times a degree-$n$ monomial $V$ in the $v_{ij}$.
The remaining expression is $b/V\in B_{S_0}$, proving membership
in the right side. If $s<n$, remove all $s$ such $Q$ factors,
$n-s$ free $q_j$ factors and $n$ free $x_j$ factors. There are
enough free $q_j$ factors since $a_q+n-s\geq n-s$. The remaining
expression is $b/(t^{n-s}V)\in B_{S_0}$, with $V$ a degree-$s$
$v$-monomial, again proving membership. Conversely,
\[
 tE_j=q_jx_j,\qquad v_{ij}E_j=x_jQ_{ij}
\]
belong to $B_{S_0}$, so every product of $n$ of these ideal
generators belongs to the left side. Multiplication by $E_j^n$
preserves distinct Laurent monomials, making membership termwise.
This proves the full ideal equality, including $\alpha=0$, when
the right side is $(t^n)$.

The complete relative kernel is
\begin{equation}\label{nb:relative-kernel}
 \ker\bigl(B_{S_0}[Y_j]\longrightarrow B_{S_0}[E_j]\bigr)
 =(tY_j-q_jx_j,\ v_{ij}Y_j-x_jQ_{ij}:1\leq i\leq\alpha).
\end{equation}
Group polynomial terms by equal Laurent images. Every kernel element
is a linear combination of binomials $bY_j^n-cY_j^m$. Cancelling
a common $Y_j$ power reduces to $bY_j^e-c$, where
$c=bE_j^e\in B_{S_0}$. By~\eqref{nb:denominator-ideals}, some
product $g_1\cdots g_e$ with each $g_l\in\{t,v_{1j},\ldots,
v_{\alpha j}\}$ divides the monomial $b$ in $B_{S_0}$.
Write $b=b'g_1\cdots g_e$ and $h_l=g_lE_j$. Each $h_l$ is
$q_jx_j$ or $x_jQ_{ij}$. The product telescoping identity expresses
\[
 b'\bigl((g_1\cdots g_e)Y_j^e-h_1\cdots h_e\bigr)
\]
as a combination of $g_lY_j-h_l$. Multiplying by the cancelled
power proves~\eqref{nb:relative-kernel} before localization,
including higher powers and torsion relations.

Adjoin the $E_j$ in order. At each stage the displayed relative
kernel generators have coefficients in the original $B$. Their
preimage in the preceding polynomial presentation adds precisely
these generators to the known kernel. Induction gives
\begin{equation}\label{nb:full-kernel}
 \ker\bigl(B[Y_1,\ldots,Y_k]\longrightarrow A\bigr)
 =(tY_j-q_jx_j,\ v_{ij}Y_j-x_jQ_{ij}:\text{all }i,j).
\end{equation}
Thus cross-relations among several $Y_j$ are included; relations
internal to $B$ are already identities in the coefficient algebra.

Give every $Y_j$ weight one in addition to the $q$-grading on $B$.
The generators in~\eqref{nb:full-kernel} have weight one, while
substitution of $G_j$ gives
\begin{equation}\label{nb:reduction-remainders}
 tG_j-q_jx_j=pz_j,\qquad v_{ij}G_j-x_jQ_{ij}=H_{ij},
\end{equation}
whose right sides have weight zero. Represent $f\in R$ by
$F\in B[Y_1,\ldots,Y_k]$ and let $d$ be its largest weight.
If its weight-$d$ component survives substitution of $E$, it is
the leading part of $f$ and lies in $A$. If it cancels, it belongs
to~\eqref{nb:full-kernel}. Express it using homogeneous coefficients
of weight $d-1$ times the weight-one kernel generators, and use
\eqref{nb:reduction-remainders} after substitution of $G$. This
represents the same $f$ with strictly smaller largest weight.
Nonnegative integer weights ensure termination. Every leading part
therefore lies in $A$; conversely its generators are leading parts
of the corresponding generators of $B$ or $G_j$. This proves
\eqref{nb:star-graded}. The same reduction shows that any $f$ of
degree at most $d$ has a representation of weighted degree at most
$d$ in the enlarged family including the weight-zero $H_{ij}$.

We now prove saturation of the monomial algebra $A$. Write its
candidate exponent vector as $t$ exponent $K$ and nonnegative
resource exponents $a_p,a_{q_j},a_{x_j},a_{z_j},a_{u_h},a_{v_{ij}}$.
The negative-$t$ generators $P_h,E_j,Q_{ij},H_{ij}$ can have
nonnegative multiplicities $P'_h,E'_j,Q'_{ij},H'_{ij}$ precisely
subject to the resource bounds
\begin{equation}\label{nb:flow-bounds}
 \begin{gathered}
 \sum_hP'_h+\sum_{i,j}H'_{ij}\leq a_p,\qquad
 E'_j+\sum_iQ'_{ij}\leq a_{q_j},\qquad E'_j\leq a_{x_j},\\
 \sum_iH'_{ij}\leq a_{z_j},\qquad P'_h\leq a_{u_h},\qquad
 Q'_{ij}+H'_{ij}\leq a_{v_{ij}}.
 \end{gathered}
\end{equation}
If $N'$ is their total multiplicity, the required $t$ exponent is
possible exactly when $N'\geq-K$, using $K+N'$ free $t$ factors.
Unused resources are supplied by free variables.

The bounds are a single-commodity flow network. Its source has
edges to $p$ and $q_j$ of capacities $a_p,a_{q_j}$. There are
unrestricted edges $p\to u_h$, followed by $u_h$-to-sink edges of
capacity $a_{u_h}$. An edge $p\to\mathrm{gate}_j$ has capacity
$a_{z_j}$, followed by unrestricted edges to $v_{ij}$. There are
unrestricted edges $q_j\to x_j$ and $q_j\to v_{ij}$; their edges
to the sink have capacities $a_{x_j}$ and $a_{v_{ij}}$, respectively.
The four path types are $P,H,E,Q$. Every unrestricted edge may be
given finite capacity $a_p+\sum_j a_{q_j}$. For integer resource
data all capacities are integral, so integral max-flow supplies
integer multiplicities whenever real multiplicities reach the
integer demand $-K>0$. Deleting unit paths permits any smaller
required nonnegative integer flow. If $K\geq0$, free variables
already give a representation.

The exponent cone is exactly this real-feasible region: a real
generator expression gives a packing, and a packing with unused
resources gives a real generator expression. Every integer point
of the cone consequently has an integer generator expression. The
lattice is the full integer lattice because all free variables are
present. Thus the exponent monoid is its cone intersected with the
lattice. Intersecting the Laurent polynomial ring with the monomial
valuation rings for the rational cone facets, as in
Proposition~\ref{nb:monomial}, proves that $A$ is normal. This
argument includes all empty index sets.

Let $F_dR$ denote elements of total $q$ degree at most $d$, and form
the weighted Rees algebra
\[
 \mathcal E=\bigoplus_{d\geq0}F_dR\,\xi^d.
\]
The terminating reduction proves finite generation: generators are
$\xi$, the weight-zero elements $t,p,x,z,u,v,P,H$, and
$q_j\xi,Q_{ij}\xi,G_j\xi$. Including $H_{ij}$ explicitly is
essential, since its cancellation expression~\eqref{nb:H} has
larger weight before cancellation. The grading is nonnegative, and
\[
 \mathcal E/(\xi)=\operatorname{gr}R=A,\qquad
 \mathcal E[\xi^{-1}]=R[\xi,\xi^{-1}].
\]
Corollary~\ref{nb:graded-deformation} makes $\mathcal E$ normal.
Localizing gives a normal $R[\xi,\xi^{-1}]$. An element of
$\Frac(R)$ integral over $R$ belongs to that Laurent polynomial
ring and, being independent of $\xi$, belongs to $R$. Thus $R$
is normal.

Finally, inverting $t$ eliminates all fraction generators, so the
fraction field is that of its $2+3k+\beta+\alpha k$ independent
base variables. For the homogeneous origin, now give \emph{all}
base variables degree one. The original $G,P,Q$ also have degree
one, while $H$ has degree two. The degree-zero part is $\C$.
There is no linear relation among the original generators: free
variables and $P,Q$ are distinct Laurent monomials, and the monomial
$pz_j/t$ occurs in the unique generator $G_j$. Its coefficient
first forces each $G_j$ coefficient to vanish; the remaining
monomial coefficients then vanish. Hence the original degree-one
generators form a basis of the cotangent space. Their number is
$2+3k+\beta+\alpha k+(k+\beta+\alpha k)=N$. The excess is
$k+\beta+\alpha k$, positive unless $k=\beta=0$. In that case
$R=\C[t,p]$, independently of $\alpha$, and is smooth.
\end{proof}

\subsection{Identification of the entire fibres and their local dimensions}
\label{nb:completion}

\begin{proof}[Proof of Theorems~\ref{nb:D}, \ref{nb:E1} and~\ref{nb:E2}]
Lemma~\ref{nb:configuration} gives precisely the stated graph
configurations, six of them smooth. Lemma~\ref{nb:charts} gives
the exact local models at every dependent configuration, and
Propositions~\ref{nb:monomial} and~\ref{nb:star} prove them normal.
Normality descends through the factor-vector torus torsor and the
free ordered-configuration \'etale quotient. Thus $G$ is normal at
every point of its finite fibre over $L$.

The non-quasi-finite locus is closed and misses that fibre; remove
its proper image in $H$. The nonnormal locus is likewise closed
for a variety over $\C$, and its proper image also misses $L$.
After these shrinkings, $G\to H$ is finite with normal integral
source and is birational. It is therefore the entire normalization:
on affine rings, a finite normal extension with the same fraction
field equals the integral closure. In particular there is precisely
one normalization point at each of the $r+1$ configurations, or
ten configurations when $r=9$ in E2.

We also check unramifiedness scheme-theoretically. The projective
point incidence $Z_H=\{(K,p):p\in\PP(K)\cap S\}$ has finite
reduced fibre at $L$. Properness makes it finite near $L$. Its fibre
cotangent spaces vanish; base change for relative differentials and
Nakayama's lemma give unramifiedness at the fibre. Removing the
closed image of the remaining non-unramified locus gives a finite
unramified incidence nearby. In its $r$-fold fibre product the
distinct-point locus is open and closed, since diagonals of a
separated unramified morphism are open and closed. This locus is
therefore finite and unramified over $H$.

The graph with its points ordered is a closed subscheme of that
incidence near the fibre. The proper image of the locus of nonreduced
Hilbert-scheme configurations may first be removed, since it misses
the fibre. Closed immersions are unramified, and descent through the
free \'etale permutation quotient proves unramifiedness of $G\to H$.
Over $\C$, its finite unramified fibres are reduced. When the models
in D are smooth, all points over $L$ are smooth; removing the proper
image of the nonsmooth locus gives a neighbourhood whose entire
normalization source is smooth. We have used only the incidence and
bundle constructions underlying~\ref{norm:criterion}; its
all-subsets-independent hypothesis is not imposed or used at these
dependent configurations.

On the rank-$r$ locus of the third flattening, $X_r(a,b,c)$ is the
bundle of surjections $(\C^c)^*\to K$ over $K\in H$. For an
independent rank-one span, every such surjection is an $r$-term
tensor sum, and conversely a rank-$r$-flattening $r$-term sum gives
such a span. Taking reduced closures preserves this equality: a
local trivialization is a product, and the bundle projection is
smooth and open. Pulling the normalization of $H$ back along this
bundle gives a finite birational source that is an open vector
bundle over a normal variety, hence normal. It is the full
normalization of this tensor locus. Base change preserves
unramifiedness and the indicated fibres; the bundle adds only $rc$
smooth parameters. Thus every point has dimension
\[
 \dim H+rc=r(a+b+c-2).
\]

Smooth and \'etale covers preserve tangent excess. For D,
Proposition~\ref{nb:monomial} gives $(r-6)(b-3)$, with zero exactly
when $r=6$ or $b=3$. For the mixed cases, the excesses are
\[
 \begin{array}{c|c}
 \text{configuration}&\text{tangent excess}\\ \hline
 \text{E1, omitted row}&(r-6)(a-3)\\
 \text{E1, omitted column}&(r-7)(a-3)+b-r+4\\
 \text{E2, omitted row}&2a+b-11\\
 \text{E2, omitted column}&a+2b-11.
 \end{array}
\]
The first is positive since $r\geq7,a\geq4$; the second is
positive since $b-r+4\geq1$. The last two are at least five
because $a,b\geq5$ and $a+b\geq11$. These are normal singular
points, including the $k=0$ and $\beta=0$ boundary cases allowed
by Proposition~\ref{nb:star}. Adding the dimension gives exactly
the claimed embedding dimensions, and the six actual-decomposition
points remain smooth. Factor permutations transport the variety,
normalization and witness, proving the remaining assertions.
\end{proof}
